\section{Introduction}

Entities or objects which have equivalent terms or expressions
in different languages may have subtle differences in 
their use scenario or perceptions in the respective culture.
For example, dragon and loong(龙) refer to the similar legendary
reptile-like creature that can fly and exhale flames, 
in both English and Chinese. While the two concepts have influenced
each other in the past, they have quite notable cultural differences. 
In the English or the larger European culture, 
this snake-like, winged being has a generally violent
and fearsome image, which is often associated with negative 
contexts such as dungeon or hell. In the Chinese setting, however,
loong is a symbol of royal power, majesty and auspiciousness. Sometimes,
it even carries the notion of playfulness, as loong is often shown
playing with a ball. In fact, these differences in perception 
can be observed by visual distinction in image search results.
\figref{fig:dragon} and \figref{fig:loong} show the top images 
returned from Bing image
search of ``dragon'' and ``龙''\footnote{\url{http://global.bing.com/images/}.}. One can see that the popular images,
which reflect the perception of these concepts in their respective
cultures, are quite different. 

\begin{figure}[th]
\centering
\includegraphics[width=0.49\columnwidth]{img/dragon1}
\hfill
\includegraphics[width=0.49\columnwidth]{img/dragon2}
%
\caption{Images of ``dragon'' on Bing.}
\label{fig:dragon}
\end{figure}

\begin{figure}[th]
\centering
\includegraphics[width=0.49\columnwidth]{img/loong1}
\hfill
\includegraphics[width=0.49\columnwidth]{img/loong2}
\caption{Images of ``龙'' on Bing.}
\label{fig:loong}
\end{figure}

Understanding such cultural differences in the expressions 
between languages can be useful in many applications, especially
when there is insufficient contextual information. 
In machine translation of short texts (such as movie subtitles),
we want to avoid translated terms that are 
irrelevant or even offensive in another culture. 
In cross-lingual recommendation systems, we may 
only recommend products or services most relevant to 
the people of a particular culture. Finally, in
foreign language education, it is useful to know that bowls 
are made to contain rice in China while they are typically much bigger 
in Europe to hold salad. 

Previous efforts have focused on translating or paraphrasing
concepts from one language to other languages (e.g., BabelNet~\cite{Navigli:2012dn}
and YAGO3~\cite{mahdisoltani2014yago3}). 
Other examples include mining synonymous word pairs across
different languages\cite{Mikolov:2013tp}, 
or bilingual lexicon\cite{linard2015bilingual}, for translation.
%\KZ{Are these suitable exmaples, and what else?}
In other words, they pay much attention on the similarity rather
than differences between the cultures.

In this paper, we propose to develop metrics
that quantatively measure the cultural differences between equivalent
expressions of objects or entities in English and Chinese. 
Our proposed framework is language independent, such that, 
given suitable corpora, the same methodology developed here can be
applied to other languages and cultures. Our goal is to compute a globally
comparable score for any pair of English-Chinese translations, where
the English term refers to an entity. We choose to restrict our
study on entities, 
 because these are traditionally thought to 
carry identical meaning across any culture. 
Furthermore, cultural 
difference and, for the same matter, culture similarity, for 
entities can be more easily verified by human judges using the image search described above. The main 
contributions of this paper can be summarized as follows.

\begin{itemize}
\item We propose four algorithms, all based on the skip-gram word 
embedding model trained from large literature corpora, 
to calculate the cultural difference between a pair
of English-Chinese translated terms;
\item We propose an innovative evaluation method that uses image search
results to identify entity pairs which have cultural differences;
\item All four algorithms above perform reasonably well, while the
KNN-cosine method outperforms the rest and achieves 0.9 precision score for
the top 10 culturally different term-pairs;
\item We build a web-based demonstration system that allows users to
search an English or a Chinese term, and returns the cultural difference
score, a percentile ranking among all indexed terms, as well as a
visual explaination why a pair of terms is cultural different (or 
indifferent).
\end{itemize}

Next we present our approaches (\secref{sec:approach}), 
show some key evaluation results (\secref{sec:eval}),
propose a demonstration plan (\secref{sec:demo}), and 
discuss some previous related work (\secref{sec:related}).

